\SourceLeaf{14}
\begin{center}
{\Large\textsc{DISQUISITIONES ARITHMETICAE.}}

\BookRule{\linewidth}

\textsc{SECTIO PRIMA}

{\scriptsize\textsc{DE}}

\textsc{NUMERORUM CONGRUENTIA IN GENERE.}

\BookRule{0.10\linewidth}
\end{center}

\begin{center}
\emph{Numeri congrui, moduli, residua et nonresidua.}
\end{center}

\Article{1}

Si numerus $a$ numerorum $b, c$ differentiam metitur. $b$ et $c$ \emph{secundum $a$ congrui} dicuntur, sin minus, \emph{incongrui}: ipsum $a$ \emph{modulum} appellamus. Uterque numerorum $b, c$ priori in casu alterius \emph{residuum}, in posteriori vero \emph{nonresiduum} vocatur.

Hae notiones de omnibus numeris integris tam positivis quam negativis\HistoricalNote{*)}{Modulus manifesto semper \emph{absolute} i. e. sine omni signo est sumendus.} valent, neque vero ad fractos sunt extendendae. \emph{E. g.} $-9$ et $+16$ secundum modulum $5$ sunt congrui; $-7$ ipsius $+15$ secundum modulum $11$ residuum, secundum modulum $3$ vero nonresiduum. Ceterum quoniam cifram numerus quisque metitur, omnis numerus tamquam sibi ipsi congruus secundum modulum quemcunque est spectandus.

\Article{2}

Omnia numeri dati $a$ residua secundum modulum $m$ sub formula $a+km$ comprehenduntur, designante $k$ numerum integrum indeterminatum. Propositionum quas post trademus faciliores nullo negotio hinc demonstrari possunt: sed istarum quidem veritatem aeque facile quivis intuendo poterit perspicere.

\SourceLeaf{15}
Numerorum congruentiam hoc signo, $\equiv$, in posterum denotabimus, modulum ubi opus erit in clausulis adiungentes, $-16\equiv 9\ (\mathrm{mod.}\ 5)$. $-7\equiv 15\ (\mathrm{mod.}\ 11)$\HistoricalNote{*)}{Hoc signum propter magnam analogiam quae inter aequalitatem atque congruentiam invenitur adoptavimus. Ob eandem caussam ill. Le Gendre in comment. infra saepius laudanda ipsum aequalitatis signum pro congruentia retinuit, quod nos ne ambiguitas oriatur imitari dubitavimus.}.

\Article{3}

\textsc{Theorema.} \emph{Propositis $m$ numeris integris successivis}
\[
a,\ a+1,\ a+2\ \GaussFourDots\ a+m-1
\]
\emph{alioque $A$, illorum aliquis huic secundum modulum $m$ congruus erit, et quidem unicus tantum.}

Si enim $\frac{a-A}{m}$ integer, erit $a\equiv A$, sin fractus, sit integer proxime maior, (aut quando est negativus, proxime \emph{minor}, si ad signum non respiciatur) $=k$, cadetque $A+km$ inter $a$ et $a+m$, quare erit numerus quaesitus. Et manifestum est omnes quotientes $\frac{a-A}{m}$, $\frac{a+1-A}{m}$, $\frac{a+2-A}{m}$ etc. inter $k-1$ et $k+1$ sitos esse; quare plures quam unus integri esse nequeunt.

\begin{center}\emph{Residua minima.}\end{center}
\Article{4}

Quisque igitur numerus residuum habebit tum in hac serie, $0, 1, 2,\ldots m-1$, tum in hac, $0, -1, -2,\ \GaussFourDots\ -(m-1)$, quae \emph{residua minima} dicemus, patetque, nisi $0$ fuerit residuum, bina semper dari, \emph{positivum} alterum, alterum \emph{negativum}. Quae si magnitudine sunt inaequalia, alterum erit $<\frac{m}{2}$, sin secus utrumque $=\frac{m}{2}$, signi respectu non habito. Unde patet, quemvis numerum residuum habere moduli semissem non superans quod \emph{absolute minimum} vocabitur.

\emph{E. g.} $-13$ secundum modulum $5$, habet residuum minimum positivum $2$, quod simul est absolute minimum, $-3$ vero residuum minimum negativum; $+5$ secundum modulum $7$ sui ipsius est residuum minimum positivum, $-2$ negativum, simulque absolute minimum.

\begin{center}\emph{Propositiones elementares de congruentiis.}\end{center}
\Article{5}

His notionibus stabilitis eas numerorum congruorum proprietates quae prima fronte se offerunt colligamus.

\SourceLeaf{16}
\emph{Qui numeri secundum modulum compositum sunt congrui, etiam secundum quemvis eius divisorem congrui.}

\emph{Si plures numeri eidem numero secundum eundem modulum sunt congrui, inter se erunt congrui} (secundum eundem modulum).

Haec modulorum identitas etiam in sequentibus est subintelligenda.

\emph{Numeri congrui residua minima habent eadem, incongrui diversa.}

\Article{6}

\emph{Si habentur quotcunque numeri $A, B, C$ etc. totidemque alii $a, b, c$ etc. illis secundum modulum quemcunque congrui}
\[
A\equiv a,\ B\equiv b\ \text{\itshape etc.,\quad erit}\quad A+B+C+\text{\itshape etc.}\equiv a+b+c+\text{\itshape etc.}
\]

\emph{Si $A\equiv a, B\equiv b$, erit $A-B\equiv a-b$.}

\Article{7}

\emph{Si $A\equiv a$, erit quoque $kA\equiv ka$.}

Si $k$ numerus positivus, hoc est tantummodo casus particularis propos. art. praec., ponendo ibi $A=B=C$ etc., $a=b=c$ etc. Si $k$ negativus, erit $-k$ positivus, adeoque $-kA\equiv -ka$, unde $kA\equiv ka$.

\emph{Si $A\equiv a, B\equiv b$, erit $AB\equiv ab$.} Namque $AB\equiv Ab\equiv ba$.

\Article{8}

\emph{Si habentur quotcunque numeri $A, B, C$ etc. totidemque alii $a, b, c$ etc. his congrui, $A\equiv a, B\equiv b$ etc., producta ex utrisque erunt congrua, $ABC\ \text{\itshape etc.}\equiv abc\ \text{\itshape etc.}$}

Ex artic. praec. $AB\equiv ab$, et ob eandem rationem $ABC\equiv abc$; eodemque modo quotcunque alii factores accedere possunt.

Si omnes numeri $A, B, C$ etc. aequales assumuntur, nec non respondentes $a, b, c$ etc., habetur hoc theorema: \emph{Si $A\equiv a$ et $k$ integer positivus, erit $A^{k}\equiv a^{k}$.}

\Article{9}

\emph{Sit $\mathrm{X}$ functio algebraica indeterminatae $x$, huius formae}
\[
Ax^{a}+Bx^{b}+Cx^{c}+\text{\itshape etc.}
\]
\emph{designantibus $A, B, C$ etc. numeros integros quoscunque; $a, b, c$ etc. vero integros non negativos. Tum si indeterminatae $x$ valores secundum modulum quemcunque congrui tribuuntur, valores functionis $\mathrm{X}$ inde prodeuntes congrui erunt.}

\SourceLeaf{17}
Sint $f, g$ valores congrui ipsius $x$. Tum ex art. praec. $f^{a}\equiv g^{a}$ et $Af^{a}\equiv Ag^{a}$, eodemque modo $Bf^{b}\equiv Bg^{b}$ etc. Hinc
\[
Af^{a}+Bf^{b}+Cf^{c}+\text{etc.}\equiv Ag^{a}+Bg^{b}+Cg^{c}+\text{etc.}\quad\text{\itshape Q. E. D.}
\]

Ceterum facile intelligitur, quomodo hoc theorema ad functiones plurium indeterminatarum extendi possit.

\Article{10}

Quodsi igitur pro $x$ omnes numeri integri consecutivi substituuntur, valoresque functionis $\mathrm{X}$ ad residua minima reducuntur, haec seriem constituent, in qua post intervallum $m$ terminorum (designante $m$ modulum) iidem termini iterum recurrunt; sive haec series ex \emph{periodo} $m$ terminorum infinities repetita, erit formata. Sit e. g. $\mathrm{X}=x^{3}-8x+6$ et $m=5$; tum pro $x=0, 1, 2, 3$ etc., valores ipsius $\mathrm{X}$ haec residua minima positiva suppeditant, $1, 4, 3, 4, 3, 1, 4$ etc., ubi quina priora $1, 4, 3, 4, 3$ in infinitum repetuntur; atque si series retro continuatur, i. e. ipsi $x$ valores negativi tribuuntur, eadem periodus ordine terminorum inverso prodit: unde manifestum est, terminos alios quam qui hanc periodum constituant in tota serie locum habere non posse.

\Article{11}

In hoc igitur exemplo $\mathrm{X}$ neque $\equiv 0$, neque $\equiv 2$ (mod. 5) fieri potest, multoque minus $=0$, aut $=2$. Unde sequitur, aequationes $x^{3}-8x+6=0$, et $x^{3}-8x+4=0$ per numeros integros et proin, uti notum est, per numeros rationales solvi non posse. Generaliter perspicuum est, aequationem $\mathrm{X}=0$, quando $\mathrm{X}$ functio incognitae $x$, huius formae
\[
x^{n}+Ax^{n-1}+Bx^{n-2}+\text{etc.}+N
\]
$A, B, C$ etc. integri, atque $n$ integer positivus, (ad quam formam omnes aequationes algebraicas reduci posse constat) radicem rationalem nullam habere, si congruentiae $\mathrm{X}\equiv 0$ secundum ullum modulum satisfieri nequeat. Sed hoc criterium, quod hic sponte se nobis obtulit, in Sect. VIII fusius pertractabitur. Foterit certe ex hoc specimine notiuncula qualiscunque de harum investigationum utilitate efformari.

\SourceLeaf{18}
\begin{center}\emph{Quaedam applicationes.}\end{center}
\Article{12}

Theorematibus in hoc capite traditis complura quae in arithmeticis doceri solent innituntur, e. g. regulae ad explorandam divisibilitatem numeri propositi per 9, 11 aut alios numeros. \emph{Secundum modulum} 9 omnes numeri 10 potestates unitati sunt congruae: quare si numerus propositus habet formam $a+10b+100c+\text{etc.}$, idem residuum minimum secundum modulum 9 dabit, quod $a+b+c+\text{etc.}$ Hinc manifestum est, si figurae singulae numeri decadice expressi sine respectu loci quem occupant addantur, summam hanc numerumque propositum eadem residua minima praebere, adeoque hunc per 9 dividi posse, si illa per 9 sit divisibilis, et contra. Idem etiam de divisore 3 tenendum. Quoniam \emph{secundum modulum} 11, $100\equiv 1$ erit generaliter $10^{2k}\equiv 1$. $10^{2k+1}\equiv 10\equiv -1$, et numerus formae $a+10b+100c+\text{etc.}$ secundum modulum 11 idem residuum minimum dabit quod $a-b+c$ etc.; unde regula nota protinus derivatur. Ex eodem principio omnia similia praecepta facile deducuntur.

Nec minus ex praecedentibus petenda est ratio regularum, quae ad verificationem operationum arithmeticarum vulgo commendantur. Scilicet si ex numeris datis alii per additionem, subtractionem, multiplicationem aut elevationem ad potestates sunt deducendi: substituuntur datorum loco residua ipsorum minima secundum modulum arbitrarium (vulgo 9 aut 11, quoniam in nostro systemate decadico secundum hos, uti modo ostendimus, residua tam facile possunt inveniri). Numeri hinc oriundi illis, qui ex numeris propositis deducti fuerunt, congrui esse debent; quod nisi eveniat, vitium in calculum irrepsisse concluditur.

Sed quum haec hisque similia abunde sint nota, diutius iis immorari superfluum foret.

\par\bigskip\bigskip\BookRule{0.29\linewidth}
